% Source: https://artofproblemsolving.com/community/c3h1889199_number_theory
% Stable identity: aops-topic-1889199-q01
Let $d=n^2+n+1$. Since $n^3-1=(n-1)(n^2+n+1)$, we have $n^3\equiv1\pmod d$. As $2013=3\cdot671$,
\[
n^{2013}=\left(n^3\right)^{671}\equiv1\pmod d,
\]
so $d$ divides $n^{2013}+61$ if and only if $d$ divides $1+61=62$.

The number $d=n(n+1)+1$ is odd, because $n(n+1)$ is a product of consecutive integers and hence even, and it is positive, because $d=\left(n+\dfrac{1}{2}\right)^2+\dfrac{3}{4}$. The odd positive divisors of $62=2\cdot31$ are $1$ and $31$.
\begin{itemize}
\item $d=1$: $n^2+n=0$, so $n=0$ or $n=-1$.
\item $d=31$: $n^2+n-30=0$, i.e.\ $(n-5)(n+6)=0$, so $n=5$ or $n=-6$.
\end{itemize}
Therefore $n\in\boxed{\{-6,\,-1,\,0,\,5\}}$. (For instance, for $n=5$: $5^3=125=4\cdot31+1$, so $5^{2013}\equiv1$ and $5^{2013}+61\equiv62\equiv0\pmod{31}$.)
